\subsection{Sprint 3}
Nesta etapa, me dediquei a uma task específica idealizada pelo nosso AGES IV Arthur Henz. Desenvolvi uma rota GET que retorna os 
todos os resultados emocionais dos usuários filtrados por algum mês específico. Essa tarefa me exigiu maior atenção para garantir que 
os dados fossem processados e entregues de forma clara, algo que nosso ORM da API dificultou um pouco. 

Além disso, me envolvi na tentativa do nosso AGES I Luiz Filho de criar a documentação das rotas da API utilizando Swagger 
para facilitar o trabalho do frontend. Contudo, o desenvolvimento dessa \textit{feature} no NestJS acabou sendo mais 
desafiador do que o esperado, consumindo um tempo significativo e impedindo a conclusão completa 
dessa tarefa dentro do prazo.

A sprint, que já era curta, apresentou dificuldades adicionais devido a problemas no planejamento 
e no respeito ao \textit{code freeze}. Diversos merge requests foram abertos no último dia, o que gerou acúmulo 
de trabalho e aumentou o débito técnico para as próximas sprints. 
\\
\\
\fbox{%
    \parbox{1\linewidth}{%
        \setlength{\parindent}{15pt}
            \itshape O padrão deste projeto foi o seguinte: a implementação da API seguia sempre tranquila, as únicas features que não conseguimos implementar eram extras para a cliente.
            Já o desenvolvimento do frontend sempre desrespeitava o code freeze e eram abertos MRs no dia anterior a entregas.
            
            Uma lição que vou levar adiante e que ainda sinto que fico devendo, é de oferecer mais ajuda para meus colegas caso sinta necessário. Por mais que eu não estava com o time de frontend,
            não me custava nada se juntar a alguns AGES I para tentar avançar mais no projeto.
    }%
}